\chapter{Introdução} \label{chap:introducao}

A previsão do comportamento de ativos financeiros a partir de informação textual é um dos problemas mais desafiadores na interseção entre a Ciência da Computação e as Finanças. Este capítulo introduz o contexto e a motivação do estudo, define de forma precisa o conceito de sentimento adotado, delimita a lacuna de pesquisa que justifica a investigação e enuncia a questão central, os objetivos e as hipóteses. Ao final, descreve-se a organização dos demais capítulos.

\section{Contextualização e Problematização}

O interesse em compreender o que move os preços das ações é antigo. Já em 1936, Keynes advertia que as oscilações dos mercados não se explicavam apenas por fundamentos econômicos, pois as decisões dos investidores seriam guiadas, em parte, por um impulso psicológico que ele chamou de espírito animal \cite{keynes_general_1936}. Essa intuição perdeu força com a consolidação das teorias racionais de precificação ao longo do século XX, mas retornou ao centro do debate quando se passou a investigar empiricamente o papel da informação na formação dos preços. O estudo pioneiro de \citeonline{cutler_what_1989}, ao perguntar de forma direta o que move os preços das ações, encontrou que os grandes acontecimentos noticiados explicavam uma fração surpreendentemente pequena das variações do mercado. Esse resultado, longe de encerrar a questão, inaugurou uma agenda de pesquisa que persiste até hoje: até que ponto o conteúdo das notícias ajuda a explicar e a antecipar o comportamento dos preços?

Durante boa parte desse percurso, as informações qualitativas, como o teor das notícias, foram tratadas como variáveis latentes, de difícil mensuração, e a pesquisa empírica concentrava-se em dados numéricos. O avanço do Processamento de Linguagem Natural alterou esse quadro ao permitir quantificar o conteúdo textual em larga escala. Duas dimensões do texto passaram a ser estudadas de forma sistemática: a legibilidade, ligada à complexidade da comunicação, e o tom, ou sentimento textual, que diz respeito à carga otimista ou pessimista das palavras empregadas \cite{garcia_sentiment_2013}. É a segunda dimensão, o sentimento, que constitui o objeto desta dissertação, e a sua definição precisa é apresentada adiante.

A compreensão da dinâmica de precificação de ativos no mercado acionário ocupa posição central na literatura financeira há décadas. A Hipótese de Mercados Eficientes (HME), formalizada por \citeonline{fama_efficient_1970}, sustenta que os preços refletem, a cada instante, toda a informação pública disponível. Em sua forma semiforte, a HME prevê que qualquer notícia nova é incorporada ao preço de maneira quase imediata, de modo que a obtenção sistemática de retornos anormais se torna improvável. Para a pesquisa computacional, essa hipótese estabelece um patamar de referência exigente: na ausência de ineficiências exploráveis, a acurácia de qualquer modelo que tente prever a direção diária de um ativo tende a permanecer próxima de 50\%, valor equivalente ao lançamento de uma moeda honesta. Esse patamar será retomado ao longo do texto como critério de avaliação realista dos resultados.

A plausibilidade empírica da eficiência plena foi contestada pelas Finanças Comportamentais. A partir da Teoria do Prospecto de \citeonline{kahneman_prospect_1979}, consolidou-se o entendimento de que os investidores operam sob racionalidade limitada e são influenciados por heurísticas e vieses cognitivos. Esses desvios sistemáticos abrem janelas de ineficiência temporária durante o processo de ajuste dos preços, intervalo em que a informação ainda não foi plenamente incorporada e no qual modelos preditivos podem, em princípio, extrair sinal. Nesse arcabouço, a imprensa especializada cumpre o papel de principal canal de atualização das expectativas. O trabalho de \citeonline{tetlock_giving_2007} foi pioneiro ao quantificar, para o mercado norte-americano, a relação entre o pessimismo veiculado pela mídia e o subsequente movimento de queda dos retornos, acompanhado de elevação da volatilidade.

O mercado acionário brasileiro, operado pela B3, é particularmente permeável a esse fluxo de informação. As ações preferenciais da Petróleo Brasileiro S.A., negociadas sob o código PETR4, constituem um caso ilustrativo. Além de expostas aos choques internacionais do preço do petróleo, essas ações respondem de maneira aguda a intervenções e sinalizações políticas internas, uma vez que a Petrobras é uma empresa de capital misto sob influência do poder público. Essa dupla exposição, ao mercado global de \textit{commodities} e ao ambiente institucional doméstico, faz da PETR4 um objeto rico para investigar como o sentimento das notícias se relaciona ao comportamento dos preços e do risco \cite{silva_o_2018, calomiris_how_2019}.

A modelagem do risco, neste trabalho, apoia-se em uma ferramenta econométrica que convém apresentar desde já, em linhas gerais, para a leitura fluir nos capítulos seguintes. Séries financeiras exibem o fenômeno do agrupamento de volatilidade, em que períodos de grandes oscilações tendem a ser seguidos por novas grandes oscilações, e períodos calmos por novas calmarias. O modelo GARCH (\textit{Generalized Autoregressive Conditional Heteroskedasticity}), proposto por \citeonline{bollerslev_generalized_1986}, captura justamente esse comportamento ao estimar uma variância que varia no tempo em função dos choques recentes e da própria inércia da volatilidade passada. Neste estudo, o GARCH fornece uma medida diária de risco que serve tanto como atributo do modelo preditivo quanto como objeto de avaliação. A formalização completa é apresentada no Capítulo~\ref{chap:metodologia}.

\subsection{Definição operacional de sentimento}

O termo sentimento admite múltiplas acepções, o que torna indispensável fixar a definição empregada nesta dissertação. Adota-se uma definição operacional e quantitativa: sentimento é o escore de polaridade, no intervalo $[-1, +1]$, atribuído ao texto de uma notícia por um modelo de classificação supervisionado. Valores negativos indicam tom pessimista, valores positivos indicam tom otimista e o valor zero indica neutralidade. Essa definição é textual e mensurável, e distingue-se de conceitos correlatos como o humor coletivo de pesquisas de opinião ou o sentimento do investidor inferido de variáveis de mercado. A escolha por uma medida derivada diretamente das manchetes alinha-se ao objetivo de avaliar o conteúdo informacional do noticiário e permite a integração do sinal textual ao modelo quantitativo.

\section{Motivação e Lacuna de Pesquisa}

Do ponto de vista da Ciência da Computação, a extração de sentimento de textos financeiros em português brasileiro permanece um problema aberto. As primeiras abordagens aplicadas ao mercado brasileiro recorreram a técnicas de contagem de palavras, como o \textit{Bag-of-Words} e o TF-IDF, que representam o texto pela frequência dos termos a partir de dicionários predefinidos \cite{nizer_metodo_2011}. Tais técnicas descartam a ordem das palavras e o contexto sintático, o que compromete a interpretação de jargões e de construções ambíguas frequentes no noticiário de mercado. A introdução do Processamento de Linguagem Natural (PLN) baseado no Mecanismo de Atenção, por \citeonline{vaswani_attention_2017}, e do modelo BERT, por \citeonline{devlin_bert_2019}, deslocou o estado da arte para representações contextuais profundas, capazes de ponderar o significado de cada palavra em função das demais. Para o português, o BERTimbau \cite{souza_bertimbau_2020} disponibilizou um modelo de linguagem treinado de forma nativa no idioma. Um modelo de propósito geral, contudo, não constitui, por si só, um classificador de sentimento do domínio financeiro, pois seus dados de treinamento não pertencem majoritariamente a esse domínio. Essa limitação é superada por modelos especializados, como o FinBERT-PT-BR \cite{santos_finbert_2023}, obtido pelo ajuste fino do BERTimbau sobre um amplo corpus financeiro.

A revisão da literatura, detalhada no Capítulo~\ref{chap:referencial_teorico}, evidencia três lacunas que motivam esta pesquisa. A primeira é a escassez de estudos que apliquem modelos de sentimento especializados em finanças ao português brasileiro, em contraste com a abundância de trabalhos voltados ao inglês. A segunda é a ausência de investigações que integrem, em uma arquitetura unificada, três componentes complementares: o sentimento textual contextual, o controle econométrico do risco e a previsão direcional por algoritmos não lineares, aplicados a um ativo específico de elevada sensibilidade política, e não a índices agregados que diluem o risco individual. A terceira é de natureza metodológica e diz respeito à necessidade de assegurar a marcação temporal precisa das notícias, sem a qual a relação causal entre a chegada da informação e a reação do preço não pode ser testada de modo legítimo.

Dois aspectos da literatura recente reforçam a pertinência de investigar especificamente o mercado brasileiro. O primeiro é que o efeito da mídia sobre os preços não é uniforme ao longo do tempo. \citeonline{garcia_sentiment_2013} demonstrou que a reação do mercado ao tom das notícias se concentra nos períodos de recessão, quando os investidores se mostram mais suscetíveis a serem movidos pelo sentimento. O segundo é que esse efeito pode variar entre países. \citeonline{griffin_how_2011} levantaram a hipótese de que, em mercados emergentes, a mídia financeira seria menos influente do que em mercados desenvolvidos, o que torna empiricamente relevante verificar se o sentimento das notícias de fato afeta um ativo de um mercado emergente como o brasileiro. Essa heterogeneidade temporal sugere, ainda, que o efeito do sentimento deva ser examinado não apenas na média, mas ao longo de toda a distribuição do retorno e da volatilidade, perspectiva que esta dissertação incorpora nas análises do Capítulo~\ref{chap:resultados}.

A partir dessas lacunas, formula-se a questão central que orienta o trabalho: de que forma a fusão do sentimento de notícias financeiras, extraído por um modelo \textit{Transformer} especializado em finanças, com a modelagem econométrica do risco pode contribuir para a previsão da direção e da volatilidade do ativo PETR4 na B3?

\section{Objetivos}

\subsection{Objetivo Geral}
Investigar o impacto do sentimento extraído de notícias financeiras, processado por um modelo \textit{Transformer} especializado no domínio financeiro em português, sobre a previsão da direção do preço e da volatilidade do ativo PETR4, por meio de uma arquitetura de fusão de dados com algoritmos de aprendizado de máquina avaliados sob protocolo cronológico que evita o vazamento de informação.

\subsection{Objetivos Específicos}
Cada objetivo específico corresponde a uma etapa necessária do encadeamento metodológico, e a sua razão de ser é explicitada a seguir, conforme recomendado pela banca de qualificação.

\begin{itemize}
    \item Coletar um corpus de notícias financeiras dotado de marcação temporal precisa. Sem a data e a hora exatas de cada publicação, não é possível alinhar de forma causal a notícia ao pregão correspondente, o que inviabiliza qualquer teste da relação de antecedência entre informação e preço. A coleta por meio de interfaces de programação de múltiplos portais atende a esse requisito e, ao diversificar as fontes, reduz o viés associado a um único veículo.
    \item Organizar o corpus segundo uma taxonomia temática supervisionada. A definição prévia de categorias interpretáveis, como notícias da própria empresa, do mercado de petróleo ou de natureza geopolítica, permite analisar separadamente a contribuição de cada vetor de informação e substitui a modelagem não supervisionada de tópicos por uma estrutura auditável e reprodutível.
    \item Extrair o sentimento das notícias por meio de um modelo especializado em finanças. Um classificador ajustado ao domínio reduz o risco de interpretação equivocada de jargões e do sentido econômico do texto, em comparação com modelos genéricos de propósito geral.
    \item Modelar a volatilidade condicional do ativo por um modelo GARCH. O agrupamento de volatilidade é uma regularidade empírica bem documentada das séries financeiras, e a sua modelagem fornece uma medida diária de risco que cumpre dupla função, como atributo preditivo e como objeto de avaliação.
    \item Treinar e avaliar classificadores não lineares sob fusão de dados e divisão cronológica. A previsão direcional requer modelos capazes de capturar relações não lineares, e a avaliação cronológica, com um conjunto de teste consultado uma única vez, assegura que o desempenho medido reflita a capacidade de generalização fora da amostra.
\end{itemize}

\section{Hipóteses da Pesquisa}

A pesquisa parte da premissa de que o fluxo de notícias atua como mecanismo de atualização das expectativas dos investidores. Sobre essa base, formulam-se três hipóteses, apresentadas como proposições a serem testadas, e não como conclusões antecipadas.

\begin{itemize}
    \item \textbf{H1 (Impacto Direcional).} O sentimento textual das notícias, quando fundido aos preços, melhora a previsão da direção dos log-retornos da PETR4 em relação a um modelo que utiliza apenas a informação de preços.
    \item \textbf{H2 (Antecipação de Risco).} O sentimento das notícias antecede a volatilidade condicional da PETR4, exercendo poder preditivo sobre o risco do ativo.
    \item \textbf{H3 (Viés de Negatividade).} Notícias de polaridade negativa associam-se a reações de mercado mais intensas do que notícias positivas de magnitude comparável, em consonância com a assimetria documentada no processamento de informação pelos investidores.
\end{itemize}

Os resultados relativos a cada hipótese, reportados no Capítulo~\ref{chap:resultados}, têm magnitude compatível com a quase-eficiência do mercado, e a sua interpretação é conduzida com a cautela que a HME impõe.

\section{Justificativa e Relevância}

No plano acadêmico, a dissertação contribui para a interseção entre as Finanças Comportamentais e a Ciência de Dados ao avaliar de forma empírica um modelo de linguagem especializado em finanças no português brasileiro, evitando a perda de informação que decorre da tradução de dicionários estrangeiros \cite{henrique_literature_2019}. Como a maior parte da produção quantitativa se concentra no mercado norte-americano e na língua inglesa, a aplicação dessas técnicas ao contexto da B3, e especificamente a um ativo de elevada sensibilidade política, ajuda a preencher uma lacuna geográfica e idiomática.

No plano prático, o estudo entrega uma arquitetura reprodutível que integra a volatilidade econométrica à semântica contextual de domínio para a análise de risco. Os resultados, reportados com transparência quanto à sua magnitude e à sua significância estatística, oferecem subsídios para o entendimento do papel informacional da mídia na precificação da PETR4, sem que se atribua ao método capacidade preditiva além da que foi efetivamente medida.

\section{Organização do Documento}

Os capítulos seguintes estão organizados da forma descrita a seguir. O Capítulo~\ref{chap:referencial_teorico} apresenta o referencial teórico e a Revisão Sistemática da Literatura, com a \textit{string} de busca, a janela temporal, os critérios de seleção e a discussão dos trabalhos correlatos, posicionando a originalidade do estudo. O Capítulo~\ref{chap:metodologia} distingue a metodologia, isto é, a natureza e os pressupostos da pesquisa, do método propriamente dito, descrevendo a coleta de dados com marcação temporal, a taxonomia, a extração de sentimento, a modelagem do risco e a fusão de dados, e ilustra o percurso completo com um exemplo. O Capítulo~\ref{chap:resultados} reporta e discute os resultados obtidos sobre o corpus completo, incluindo os testes de significância, a análise de causalidade e a avaliação da volatilidade. O Capítulo~\ref{chap:conclusao} sintetiza as conclusões, examina as limitações, discute a generalização para outros ativos e aponta os desdobramentos futuros.
